\section{Artifact Index of the Lean 4 Verification}
\label{app:artifact}

Section~\ref{ch:verify} reports the machine verification of the algebraic skeleton. This appendix makes that claim inspectable without access to the artifact: it records where the artifact lives, how it is reproduced, what its axiom budget is, which paper item each machine-checked declaration discharges, and under exactly which premises the load-bearing declarations hold.

\subsection{Layout, Reproduction, and Axiom Budget}

The artifact is a self-contained Lean~4 project distributed in the \texttt{lean-proof} directory of the project repository at
\begin{center}
\texttt{https://github.com/aki-xavier/frozen-axis}
\end{center}
under the \texttt{lean-proof} path. It consists of four files: \texttt{FormalProof.lean}, the single development file of 1004 lines containing every definition, lemma, proposition, theorem, the final theorem, the example variants, and the generalization upgrades; \texttt{lakefile.lean}, which pins Mathlib to the matching release; \texttt{lake-manifest.json}, which locks every transitive dependency to an exact revision; and \texttt{lean-toolchain}, which pins the toolchain to \texttt{leanprover/lean4:v4.34.0}. No component outside this directory is required.

Reproduction is a single command from inside the directory:
\begin{center}
\texttt{cd lean-proof \&\& lake build}
\end{center}
The build was executed on the developer's machine at the blob hash recorded in Section~\ref{ch:verify} and terminates with zero errors, zero warnings, and zero occurrences of \texttt{sorry}; the development also contains no \texttt{axiom} declaration of its own. The reproducible form of this claim, stronger than the \texttt{sorry}-count, is the axiom budget: the command \texttt{\#print axioms} applied to the load-bearing declarations returns
\begin{center}
\texttt{[propext, Classical.choice, Quot.sound]}
\end{center}
for each of \texttt{finale}, \texttt{lemma1\_sharp}, and \texttt{tail\_realization}, and, by the same command, for the remaining declarations of the development. These three are Lean's own standard axioms, the ones every Mathlib development inherits; in particular the list contains no \texttt{sorryAx}, so no step of the chain is discharged by an unproved assertion. The artifact is therefore auditable at the level of the proof term, not merely at the level of a reported build outcome.

One scope remark belongs here, because it delimits what the artifact does and does not certify. The verified content is the deductive chain: definitions, the mask criterion, Lemmas 1 to 3, Propositions 4 and 6, Theorems 5 and 7, and their composition in \texttt{finale}. The premises of that chain which concern the modeling side rather than the algebra are not certified, because they are not theorems: the satisfaction of Criterion P4 on real tasks appears as an explicit premise parameter, the continuous-side Euler--Lagrange derivation is not formalized, including the differentiation that produces \eqref{eq:focgeneral} from a general pairwise potential, and network nondegeneracy is a checkable property. Section~\ref{sec:layers} assigns each component its tier, and Table~\ref{tab:layers} is the authoritative statement of that boundary.

\subsection{Index from Paper Items to Machine-Checked Declarations}

Table~\ref{tab:artifact} lists the correspondence. The tier column repeats the assignment of Table~\ref{tab:layers} so that the reader can see, for each paper item, whether its machine-checked form is carrier-independent (Tier A), anchored on a Mathlib structure theorem (Tier B), or outside the machine guarantee (Tier C).

\begin{table}[htbp]
\caption{Index from paper items to Lean 4 declarations. Declaration names are those of \texttt{FormalProof.lean}. Tier A: carrier-independent algebraic core; Tier B: structured general forms anchored on Mathlib; Tier C: modeling components not covered by the machine guarantee.}
\label{tab:artifact}
\centering
\small
\renewcommand{\arraystretch}{0.88}
\begin{tabular}{@{}>{\raggedright\arraybackslash}p{0.26\textwidth}>{\raggedright\arraybackslash}p{0.56\textwidth}>{\centering\arraybackslash}p{0.09\textwidth}@{}}
\hline
\textbf{Paper item} & \textbf{Lean declaration} & \textbf{Tier}\\
\hline
Definition 1 (ill-posedness), Criterion P1 & \texttt{IllPosed}, \texttt{P1} & B\\
\hline
Mask criterion, both directions & \texttt{illPosed\_of\_hole}, \texttt{not\_illPosed\_of\_full} & B\\
\hline
Mask criterion in kernel form & \texttt{illPosed\_iff\_nontrivial\_ker} & B\\
\hline
Definition 2 (context rule) & \texttt{IndexedContextRule}, \texttt{IsScalarFamily}, \texttt{IsFrozenFamily} & A\\
\hline
Definition 3 (Mode A / Mode B) & \texttt{IsModeA}, \texttt{DependsOnlyOn}, \texttt{IsModeB} & A\\
\hline
Definition 4 (Paradigm E / I) & \texttt{ParadigmE}, \texttt{ParadigmI}, \texttt{IsParadigmI} & A\\
\hline
Criterion P2 (heterogeneity) & \texttt{P2}, \texttt{p2\_gives\_nonconst}, \texttt{hetero\_gives\_nonconst} & B\\
\hline
Criterion P3 (distributional openness; mask form) & \texttt{P3}, \texttt{p3\_implies\_p1} & A\\
\hline
Criterion P4, component form (Section~\ref{sec:p4}) & premise parameter \texttt{hGapPositive}; the barrier, bounded-description, and realization conditions and the barrier-to-gap inequality are prose only & C\\
\hline
Lemma 1 & \texttt{lemma1\_sharp}, \texttt{lemma1\_final}, \texttt{foc\_proportionality} & A\\
\hline
Lemma 1, general form, \eqref{eq:proportiongeneral} and \eqref{eq:tablecondition} & \texttt{general\ttu proportionality}, \texttt{scalar\ttu family\ttu iff\ttu ratio\ttu agrees}, \texttt{frozen\ttu family\ttu two\ttu configs\ttu iff}, \texttt{position\ttu wise\ttu escape\ttu witness} & A\\
\hline
Lemma 2, extrinsic side & \texttt{lemma2\_full}, \texttt{lemma2\_E\_stationary\_infeasible} & A\\
\hline
Lemma 2, endogenous side & \texttt{paradigmI\_unique\_minimizer}, \texttt{paradigmI\_inner\_unique}, \texttt{paradigmI\_L2\_ae} & A and B\\
\hline
Lemma 3, toy form & \texttt{lemma3\_final2} & B\\
\hline
Lemma 3, general measure form & \texttt{lemma3\_general}, \texttt{tail\_realization} & B\\
\hline
Lemma 3, gap mechanism, both demands of Example~\ref{ex:gap} & \texttt{toy\_positive\_gap}, \texttt{modeB\_positive\_gap}, \texttt{modeA\_loss\_attained}; \texttt{evidence\_gap\_is\_full}, \texttt{modeB\_evidence\_gap}, \texttt{modeA\_evidence\_attained} & A\\
\hline
Propositions 4 and 6 (NN / BN are Mode B) & \texttt{neural\_isModeB}, \texttt{neural\_not\_modeA}, \texttt{spn\_isModeB}, \texttt{spn\_not\_modeA} & A\\
\hline
Nondegeneracy lemma & \texttt{nondeg\_not\_dependsOn} & A\\
\hline
Collision construction & \texttt{collision\_of\_masked}, \texttt{collision\_general} & A and B\\
\hline
Theorems 5 and 7 (closures) & \texttt{theorem5}, \texttt{theorem7\_spn} & A\\
\hline
Final theorem & \texttt{finale} & A and B\\
\hline
Continuous-side optimality & Euler--Lagrange derivation; not formalized & C\\
\hline
Network nondegeneracy & checkable property; not formalized & C\\
\hline
\end{tabular}
\end{table}

Three entries in Table~\ref{tab:artifact} deserve emphasis because they are the ones a reader is most likely to look for. First, Criterion P4 is listed as Tier C and as a premise parameter rather than as a formalized definition: the verified form of Lemma 3 consumes \texttt{hGapPositive}, the hypothesis that the risk gap exceeds the bound, and does not attempt to formalize either the criterion or its component form of Section~\ref{sec:p4}, in which the barrier, bounded-description, and realization conditions are combined into the risk-gap floor by the barrier-to-gap inequality. This is the same dislocation already stated in Section~\ref{ch:verify}, recorded here in the index so that it cannot be missed. Second, the collision construction and the classification and closure declarations are Tier A: their proofs quantify over arbitrary index types, so the toy carrier plays no role in them and the toy end and the continuous end are both instantiations of one proof. Third, Definition~2's indexed family is listed as Tier~A rather than as Tier~C: the family, its constant and frozen specializations, and the algebra of the demand ratio are machine-checked, so what a reader cannot verify from the artifact is not the indexed reading but the step in which a general pairwise potential becomes the coefficient \eqref{eq:focgeneral}, together with the operator-valued couplings that the general form does not claim.

\subsection{Machine-Checked Statements of the Load-Bearing Declarations}

The index names the declarations; the premises are what a reader must inspect to judge the strength of the claim. The listing below gives the statements of the load-bearing declarations verbatim in their logical content, transliterated to ASCII for typesetting (\texttt{R} for the reals, \texttt{not}, \texttt{forall}, \texttt{exists}, \texttt{and}, \texttt{->}, \texttt{!=}, \texttt{>=}, \texttt{<=} for the corresponding symbols). The Unicode source is in the artifact.

The two declarations that carry the shape exclusion of Lemmas 1 and 2:

\begin{quote}
\small
\begin{verbatim}
theorem lemma1_sharp {x d : Depth} {lambda : R} {i j : Pixel}
    (h1 : FirstOrder x d lambda i) (h2 : FirstOrder x d lambda j)
    (hHetero : (x i - d i) * kappa d j != (x j - d j) * kappa d i) :
    False

theorem lemma2_full {d : Depth} {i j : Pixel}
    (hHetero : kappa d i != kappa d j) :
    (not exists lambda, lambda != 0 and
      ParadigmEFirstOrder d lambda i and ParadigmEFirstOrder d lambda j)
    and forall pred, IsParadigmI pred ->
        exists p, (forall y, pred p <= pred y) and
                  (forall y, pred y = pred p -> y = p)
\end{verbatim}
\end{quote}

The general form of Lemma 1, in which \texttt{delta} is the derivative of the separable data term at each position and the demand ratio takes the place of the residual. No convexity appears in any of these statements, because each is an implication between stationarity conditions:
\begin{quote}
\small
\begin{verbatim}
theorem general_proportionality {delta : Pixel -> R} {d : Depth} {lambda : R}
    {i j : Pixel}
    (h1 : delta i + lambda * kappa d i = 0)
    (h2 : delta j + lambda * kappa d j = 0) :
    delta i * kappa d j = delta j * kappa d i

theorem scalar_family_iff_ratio_agrees (delta : Pixel -> R) (d : Depth)
    {k : Pixel} (hk : kappa d k != 0) :
    (exists R : IndexedContextRule, IsScalarFamily R and
        forall i, GeneralFirstOrder delta d R i) <->
      forall i, delta i * kappa d k = delta k * kappa d i

theorem frozen_family_two_configs_iff (delta delta' : Pixel -> R)
    (d d' : Depth) (hkappa : forall i, kappa d i != 0) :
    (exists R : IndexedContextRule, IsFrozenFamily R and
        (forall i, GeneralFirstOrder delta d R i) and
        (forall i, GeneralFirstOrder delta' d' R i)) <->
      forall i, delta i * kappa d' i = delta' i * kappa d i

theorem frozen_family_infeasible_of_ratio_differs
    {delta delta' : Pixel -> R} {d d' : Depth} {i : Pixel}
    (hdiff : delta i * kappa d' i != delta' i * kappa d i) :
    not exists R : IndexedContextRule, IsFrozenFamily R and
        (forall j, GeneralFirstOrder delta d R j) and
        (forall j, GeneralFirstOrder delta' d' R j)

theorem position_wise_escape_witness {delta : Pixel -> R} {d : Depth}
    {i j : Pixel} (hkappa : forall k, kappa d k != 0)
    (hne : delta i * kappa d j != delta j * kappa d i) :
    (exists R : IndexedContextRule, IsFrozenFamily R and
        IsPositionWise R and forall k, GeneralFirstOrder delta d R k) and
      not exists R : IndexedContextRule, IsScalarFamily R and
        forall k, GeneralFirstOrder delta d R k
\end{verbatim}
\end{quote}

Two things about this group are worth noting, because they are where the index becomes a formal object rather than a remark. The scalar family is the special case of the frozen family with one common entry, and \texttt{scalar\ttu family\ttu iff\ttu ratio\ttu agrees} and \texttt{frozen\ttu family\ttu two\ttu configs\ttu iff} are the two halves of the corresponding obstruction: at one configuration a scalar exists exactly when the demand ratio is the same at every position, and across two configurations a frozen family exists exactly when the demand ratio at each position is the same in both. The escape witness separates the two by construction, exhibiting a family that satisfies every condition at one configuration while no scalar family does. And the classification of such a table is per position, by the same collision argument as the scalar carrier:
\begin{quote}
\small
\begin{verbatim}
theorem frozen_table_not_modeA (W : Pixel -> R) (i : Pixel) (hW : W i != 0)
    (d e : Depth) (hd0 : d 0 = e 0) (hd1 : d 1 != e 1) :
    not IsModeA (tableRule W i)
\end{verbatim}
\end{quote}
\noindent This has the shape of \texttt{neural\ttu not\ttu modeA} with the scalar entry \texttt{W} replaced by the table entry \texttt{W i}, so the number of entries does not enter: it is the machine-checked form of the sentence that one scalar is a parameter table and a million parameters are a parameter table as well.

The measure-theoretic core of Lemma 3, in its general form and with upper-tail realization demoted from premise to theorem:

\begin{quote}
\small
\begin{verbatim}
theorem tail_realization {Omega : Type} [MeasurableSpace Omega]
    {mu : Measure Omega} [IsProbabilityMeasure mu] {f : Omega -> R}
    (hf : Integrable f mu) :
    exists omega, integral f dmu <= f omega

theorem lemma3_general {Omega : Type} [MeasurableSpace Omega]
    {mu : Measure Omega} [IsProbabilityMeasure mu]
    {loss : Omega -> R} (hloss : Integrable loss mu)
    {rmin eps : R} (hmin : 0 <= rmin) (hgap : integral loss dmu - rmin > eps)
    (hbound : forall omega, loss omega <= eps) :
    False
\end{verbatim}
\end{quote}

Note what is and is not a premise here. In \texttt{lemma3\_general} the upper-tail realization premise no longer appears: it is discharged by \texttt{tail\_realization}, which holds for an arbitrary probability measure and an arbitrary integrable function. The premises that remain are the two the paper declares: a risk gap exceeding the bound, and nonnegativity of the minimum risk. The carrier side, the collision argument, and the two closures:

\begin{quote}
\small
\begin{verbatim}
theorem neural_not_modeA (W : R) (hW : W != 0)
    (d e : Depth) (hd0 : d 0 = e 0) (hd1 : d 1 != e 1) :
    not IsModeA (neuralRule W)

theorem theorem5 (sys : System) (hP3 : P3 sys) (W : R) (hW : W != 0) :
    not IsModeA (neuralRule W)

theorem theorem7_spn (sys : System) (hP3 : P3 sys)
    (p : SpnParams) (hw : p.w != 0) :
    not IsModeA (spnRule p)
\end{verbatim}
\end{quote}

Finally the composition itself, which is the existential negation's algebraic form in one conjunction:

\begin{quote}
\small
\begin{verbatim}
theorem finale (sys : System) (hP3 : P3 sys)
    (m : MeasureSetup) (eps : R) (heps : eps > 0) (W : R) (hW : W != 0)
    (hTail : exists d, m.lossR (neuralRule W) d >= m.riskP (neuralRule W))
    (hGapPositive : m.gapP (neuralRule W) > eps)
    (hNonnegMin : m.riskMin (neuralRule W) >= 0) :
    not IsModeA (neuralRule W)
    and not GeneralizationBound m eps (neuralRule W)
\end{verbatim}
\end{quote}

The premise list of \texttt{finale} is the complete inventory of what the machine-checked composition assumes: Criterion P3 as a hypothesis on the system, the three measure-theoretic premises of Lemma 3, and the nondegeneracy \texttt{W != 0} of the carrier's parameter. Nothing else enters, and in particular no premise is smuggled in as a definitional unfolding; the gap premise is assumed here directly, and its derivation from the barrier, bounded-description, and realization conditions of Section~\ref{sec:p4} is prose. This is the concrete sense in which the artifact is inspectable: the reader can check each of these premise lists against the prose of Section~\ref{ch:falsify} and see exactly where the two agree and where, in the case of P4 and in the differentiation of the general pairwise potential, the verified form is explicitly weaker.

\subsection{Maintaining the Correspondence}

The index of Table~\ref{tab:artifact} and the statements of the previous subsection are kept in correspondence with the sources by the repository's own layout: \texttt{README.md} carries the same index in Markdown, \texttt{package-lean-proof.sh} packages the sources together with the two version-lock files while excluding the build cache, and the single-file submission version of this paper is regenerated from \texttt{sections/} rather than edited by hand. A reader who finds a discrepancy between a statement above and the source should treat the source in \texttt{lean-proof/FormalProof.lean} as authoritative and report it; the index is a map, not a second copy of the development.
